\documentclass[12pt]{article}

\usepackage{amsmath}
\usepackage{booktabs}
\usepackage{caption}
\usepackage[margin=1in]{geometry}
\usepackage{natbib}

\usepackage{my_article_header}

\title{Monetary Tightening and U.S.\ Bank Fragility in 2023:\\
       Replication of Jiang et al.\ (2024)}
\author{Summer Han \and Joe Wang\\[4pt]
        \small FINM 32900 --- Final Project}
\date{March 2025}

\begin{document}

\maketitle

\begin{abstract}
We replicate the core empirical findings of \citet{jiang2024monetary},
which quantifies mark-to-market (MTM) losses on U.S.\ commercial bank
asset portfolios resulting from the 2022--2023 Federal Reserve tightening
cycle.  Using Q1~2022 WRDS Call Report balance sheet data and iShares ETF
price changes from Q1~2022 to Q1~2023, we compute asset-level MTM losses
for all U.S.\ commercial banks classified by size: Small ($\leq\$1.384$B),
Large non-GSIB, and Global Systemically Important Banks (GSIBs).
Our replication reproduces Table~1 and the supporting balance sheet
composition statistics in Table~A1 and Figure~A1.
\end{abstract}

\tableofcontents
\newpage

%=============================================================================
\section{Introduction}
%=============================================================================

The 2022--2023 Federal Reserve rate hiking cycle---the most aggressive in
four decades---caused substantial mark-to-market losses on U.S.\ bank fixed
income portfolios.  \citet{jiang2024monetary} estimate that aggregate
unrealized losses on the banking system's assets reached approximately
\$2.2~trillion by Q1~2023, raising questions about systemic fragility,
uninsured deposit exposure, and the risk of depositor runs analogous to
Silicon Valley Bank's collapse in March 2023.

This replication report reconstructs the paper's main empirical exercise
using publicly available data from WRDS (FFIEC Call Reports) and iShares
ETF prices sourced via Yahoo Finance.  Section~\ref{sec:data} describes the
data sources.  Section~\ref{sec:method} explains the MTM loss methodology.
Section~\ref{sec:results} presents our replicated tables and figure.

%=============================================================================
\section{Data}\label{sec:data}
%=============================================================================

\subsection{Call Report Balance Sheet Data}

Bank-level balance sheet data come from the FFIEC Call Reports accessed
through WRDS.  We use the Q1~2022 quarterly snapshot (report date
2022-03-31) from five WRDS tables:
\texttt{bank.wrds\_call\_rcon\_1}, \texttt{bank.wrds\_call\_rcon\_2},
\texttt{bank.wrds\_call\_rcfd\_1}, \texttt{bank.wrds\_call\_rcfd\_2},
and \texttt{bank.wrds\_call\_rcfn\_1}.
RCON codes cover domestic offices only; RCFD codes cover consolidated
domestic and foreign offices; RCFN code covers RCFN codes cover foreign offices only.

For the extended version, we obtain Single Period Call Reports directly
from the FFIEC website through a web-scraping pipeline. These data span
2022~Q4 to 2025~Q4 (with report date 2024-12-31). To maintain
compatibility with the original pipeline, the scraped data are organised
into the same five table structures used in the WRDS-based version.

The balance sheet items extracted by maturity bucket (less than 3 months,
3m--1y, 1y--3y, 3y--5y, 5y--15y, greater than 15 years) include:
\begin{itemize}
  \item Residential mortgage-backed securities (RMBS), codes
        \texttt{rcfd/rcona555--560}
  \item Treasury and non-RMBS securities, codes \texttt{rcfd/rcona549--554}
  \item First-lien residential mortgage loans, codes
        \texttt{rcfd/rcona564--569}
  \item Other loans and leases, codes \texttt{rcfd/rcona570--575}
\end{itemize}
Total assets (\texttt{rcfd/rcon2170}), uninsured deposits
(\texttt{rcon5597}), and insured deposits (\texttt{rconf049} +
\texttt{rconf045}) are also extracted.

\subsection{ETF Price Data}

Interest rate sensitivity for each maturity bucket is proxied by iShares
U.S.\ Treasury ETF price changes over the period Q1~2022 to Q1~2023:

\begin{center}
\begin{tabular}{lll}
\toprule
Maturity Bucket & ETF Ticker & Description \\
\midrule
$<$3m, 3m--1y          & SHV  & iShares Short Treasury Bond \\
1y--3y                 & SHY  & iShares 1--3 Year Treasury Bond \\
3y--5y                 & IEI  & iShares 3--7 Year Treasury Bond \\
5y--15y (70\% weight)  & IEF  & iShares 7--10 Year Treasury Bond \\
5y--15y (30\% weight)  & TLH  & iShares 10--20 Year Treasury Bond \\
$>$15y                 & TLT  & iShares 20+ Year Treasury Bond \\
MBS multiplier         & MBB  & iShares MBS ETF \\
Treasury index         & GOVT & iShares U.S.\ Treasury Bond \\
\bottomrule
\end{tabular}
\end{center}

ETF prices are downloaded via \texttt{yfinance} and resampled to quarterly
frequency. Two methodological notes on ETF proxies: (1)~the paper uses the
S\&P 3-5 Year Treasury Bond Index (SPBDUS5T) for the 3y--5y bucket, but we
use IEI as a proxy due to data availability; (2)~the WRDS \texttt{5y--15y}
maturity bucket spans bonds from 5 to 15 years, so IEF~(7--10yr) alone
understates losses for the longer-duration portion of the range.  We apply a
70/30 blended price change of IEF and TLH~(10--20yr) to better approximate
the effective duration distribution of holdings across the full bucket.

\subsection{Interest Rate Shock}

Table~\ref{tab:table_etf} reports the actual price changes experienced by
each ETF over the measurement window.  The table serves as a quantitative
summary of the interest rate shock driving all downstream loss estimates:
readers can see exactly how steeply price declines scaled with duration,
and can verify the RMBS multiplier used to adjust mortgage-related losses.

\clearpage
\input{../_output/table_etf.tex}

%=============================================================================
\section{Methodology}\label{sec:method}
%=============================================================================

\subsection{Bank Size Classification}

Banks are classified into three groups following \citet{jiang2024monetary}:
\begin{itemize}
  \item \textbf{GSIB}: 17 Global Systemically Important Banks identified by
        RSSD ID, including U.S. chartered banks affiliated with GSIB
        holding companies.
  \item \textbf{Large non-GSIB}: Total assets $>$ \$1.384 billion (threshold
        from the paper).
  \item \textbf{Small}: Total assets $\leq$ \$1.384 billion.
\end{itemize}

\subsection{Mark-to-Market Loss Calculation}

For each bank $i$ and maturity bucket $b$, the MTM loss is:
\begin{equation}
  \text{Loss}_{i,b} = H_{i,b} \times \Delta P_b
\end{equation}
where $H_{i,b}$ is the book value of holdings in bucket $b$ and $\Delta P_b$
is the proportional price change of the matching ETF from Q1~2022 to
Q1~2023.

\paragraph{RMBS multiplier.}
RMBS and first-lien mortgage holdings are more sensitive to mortgage rate
movements than to the general Treasury yield curve.  Following
\citet{jiang2024monetary}, we scale the Treasury price changes by an
RMBS multiplier:
\begin{equation}
  m = \frac{\Delta P_{\text{MBS}}}{\Delta P_{\text{Treasury Index}}}
    = \frac{P^{\text{MBB}}_{\text{end}}/P^{\text{MBB}}_{\text{start}} - 1}%
           {P^{\text{GOVT}}_{\text{end}}/P^{\text{GOVT}}_{\text{start}} - 1}
\end{equation}
RMBS losses in bucket $b$ are then $H_{i,b}^{\text{RMBS}} \times m \times
\Delta P_b$, and first-lien mortgage losses are computed identically.
Treasury/other securities and other loans use $\Delta P_b$ directly (no
multiplier).

\subsection{Uninsured Deposit Metrics}

Two solvency metrics are computed per bank:
\begin{enumerate}
  \item \textbf{Uninsured deposit ratio}: uninsured deposits as a fraction
        of total assets.
  \item \textbf{Insured deposit coverage}: insured deposits divided by
        mark-to-market asset value (total assets minus aggregate MTM loss).
\end{enumerate}

%=============================================================================
\section{Results}\label{sec:results}
%=============================================================================

Table~\ref{tab:table1} replicates Table~1 of \citet{jiang2024monetary},
reporting aggregate MTM losses and uninsured deposit ratios by bank size
group.  Table~\ref{tab:table_a1} replicates Appendix Table~A1, showing the
balance sheet composition of U.S.\ commercial banks as of Q1~2022.
Figure~\ref{fig:figure_a1} replicates Figure~A1, plotting the distribution
of asset categories across the bank size spectrum.
Figure~\ref{fig:figure_fragility} provides our own analysis: a bank-level
scatter plot showing which banks were simultaneously exposed to large MTM
losses \emph{and} high uninsured deposit ratios---the combination most
conducive to a depositor run.

\clearpage
\input{../_output/table1.tex}

\clearpage
\input{../_output/table_a1.tex}

\clearpage
\begin{figure}[htbp]
  \centering
  \includegraphics[width=\textwidth]{../_output/figure_a1.pdf}
  \caption{Balance sheet composition of U.S.\ commercial banks by size as of
           Q1~2022.  Each bar shows the fraction of total assets in each
           category.  Replicates Figure~A1 of \citet{jiang2024monetary}.}
  \label{fig:figure_a1}
\end{figure}

\clearpage
\subsection{Bank Fragility: Loss Exposure vs.\ Uninsured Deposit Risk}

Figure~\ref{fig:figure_fragility} plots the joint distribution of MTM
loss rate and uninsured deposit reliance for every bank in the sample.
The two dimensions together define a bank's fragility: a large loss alone
need not be destabilizing if depositors are fully insured, and high
uninsured deposits alone need not cause a run if asset values are intact.
It is the combination---upper-right corner of the chart---that
characterizes banks most vulnerable to the kind of self-fulfilling run
that brought down Silicon Valley Bank in March 2023.

\begin{figure}[htbp]
  \centering
  \includegraphics[width=0.88\textwidth]{../_output/figure_fragility.pdf}
  \caption{Fragility Map: Mark-to-Market Loss Rate vs.\ Uninsured Deposit
           Exposure by Bank Size.  Each point represents one of the 4,844
           U.S.\ commercial banks as of Q1~2022.  The horizontal axis shows
           each bank's MTM loss as a percentage of book assets; the vertical
           axis shows uninsured deposits as a percentage of mark-to-market
           assets.  Banks in the upper-right corner are most fragile: they
           simultaneously suffered large valuation losses and fund themselves
           primarily with uninsured depositors who may run if losses are
           recognized.  The dashed horizontal line marks the 50\% uninsured
           threshold, above which a majority of the deposit base is at risk of
           flight.  Large non-GSIB banks and GSIBs tend to have elevated
           uninsured ratios (reflecting institutional and wholesale funding
           bases), while small community banks are more dispersed.  The shaded
           region (loss $>9\%$, uninsured $>50\%$) contains the banks most
           analogous to Silicon Valley Bank's pre-failure risk profile.}
  \label{fig:figure_fragility}
\end{figure}

\clearpage
%=============================================================================
\section{Conclusion}
%=============================================================================

Our replication confirms the main findings of \citet{jiang2024monetary}:
the 2022--2023 rate hiking cycle inflicted large unrealized losses on U.S.\
bank asset portfolios, with the aggregate MTM loss equaling a substantial
fraction of the banking system's total equity.  Uninsured depositors,
particularly at mid-size banks, faced meaningful insolvency risk if MTM
losses were recognized.  These results highlight the role of interest rate
risk in bank fragility and the importance of accounting for unrealized losses
in supervisory stress tests.

%=============================================================================
% FFIEC Extension (2024-2025) — included only if outputs exist
%=============================================================================
\IfFileExists{../_output/table1_ffiec.tex}{%
\clearpage
\section{Extension: Updated Analysis Using FFIEC Data (Q4 2023--Q4 2025)}

\subsection{Motivation}

The original \citet{jiang2024monetary} analysis captures a single snapshot---Q1~2022 balance
sheets marked to Q1~2023 prices---near the beginning of the Federal Reserve's tightening
cycle.  To understand the full arc of the episode, one needs to track conditions through the
peak and into the subsequent easing.  The Fed held the federal funds rate at
5.25--5.50\% from July~2023 through September~2024, the highest level in two decades,
before cutting rates by a cumulative 100 basis points through the end of 2024.  During this
window, the banking sector operated under the implicit subsidy of the Bank Term Funding
Program (BTFP), which allowed banks to borrow against par value of Treasury and
agency securities through March~2024, partially alleviating liquidity pressure from
unrealized losses.

This extension asks: \emph{did bank fragility persist, diminish, or worsen as the rate
environment evolved?}  The analysis uses the same ETF-based MTM loss methodology applied to
balance sheet snapshots from Q4~2023 to Q4~2025, sourced directly from FFIEC Single Period
Call Reports rather than WRDS.  The ETF price change window shifts accordingly---Q4~2023
balance sheets are marked to Q4~2024 prices, and so on---capturing the incremental
loss or gain from interest rate movements within each annual holding period.

\subsection{Results}

Table~\ref{tab:table1_ffiec} replicates the Table~1 structure for the FFIEC extension
period.  The key comparison is the aggregate loss and Loss/Asset ratio relative to
the original Q1~2022 snapshot: a lower loss rate would indicate that rising rates have
been fully priced in (or that banks actively de-risked fixed-income portfolios following
the SVB episode), while a persistently elevated loss rate would suggest that system-wide
vulnerabilities remain.

\input{../_output/table1_ffiec.tex}

\IfFileExists{../_output/table_a1_ffiec.tex}{%
\clearpage

Table~\ref{tab:table_a1_ffiec} shows the balance sheet composition as of the FFIEC
extension snapshot date.  Changes in the securities share and the residential mortgage
share relative to Table~A1 (Q1~2022) reflect portfolio rebalancing decisions made in
response to the rate environment: banks that sold longer-duration securities realized
losses but reduced future mark-to-market exposure; those that held experienced continued
unrealized depreciation.

\input{../_output/table_a1_ffiec.tex}
}{}

\IfFileExists{../_output/figure_a1_ffiec.pdf}{%
\clearpage
\begin{figure}[htbp]
  \centering
  \includegraphics[width=\textwidth]{../_output/figure_a1_ffiec.pdf}
  \caption{Balance sheet composition of U.S.\ commercial banks as of the FFIEC
           extension snapshot date.  Analogous to Figure~A1 but for the extended
           sample period.  Compare with Figure~\ref{fig:figure_a1} (Q1~2022) to
           assess portfolio rebalancing over the tightening cycle.}
  \label{fig:figure_a1_ffiec}
\end{figure}

Figure~\ref{fig:figure_a1_ffiec} shows the aggregate asset and liability mix for the
extended period.  A shift away from long-duration securities relative to the Q1~2022
baseline (Figure~\ref{fig:figure_a1}) would signal active de-risking; a stable or
growing securities share would indicate that banks continued to hold impaired portfolios.
The liability side is equally informative: rising uninsured deposit ratios would suggest
that institutional depositors remain a potential run channel even as the acute crisis
of 2023 receded.

Together, these results provide a dynamic view of U.S.\ bank fragility across the full
rate cycle---from the shock origination in 2022, through the peak in 2023--2024, and
into the early easing phase of 2024--2025.
}{}
}{}% end \IfFileExists{table1_ffiec}
 

%=============================================================================
\section{Reflection}
%=============================================================================

\subsection{Joe Wang}
My favorite part of this project was using doit to automate the data processing and replication steps.  In the middle of the semester, I was working on a quant competition, where we had
to write a 10 page paper on stablecoin microstructure and have reproducable code alongside it, and I used doit to create the pipeline. It made it much easier to manage the data and 
collaboration with my teammates. In this project, it was slightly different since we had the previous year's group's code to begin with. This made it a little tricky since refactoring the 
existing code was honestly harder than starting from scratch since there were so many moving parts. I used Claude to create mappings for me to visualize folder structures and save time, which was 
very helpful and also was my first time using ClaudeCode. That was a great experience since I don't think i would have tried ClaudeCode otherwise.

The biggest challenge to me was working with the data from WRDS. The project made me realize how much work goes into processing the data since there were so many individual variables that must
be understood and extracted. Without the previous year's code, I feel like I would have been completely lost. 


\subsection{Summer Han}
My favourite part of this project was how I ended up learning about bank fragility and mark-to-market losses not from a textbook, 
but backwards, through last year’s group’s code and data. I had no background in 
fixed income or banking regulation going in. But tracing through the pipeline, 
working out what each variable actually represents, and understanding why RMBS 
requires a different multiplier from Treasuries,  that process of reverse-engineering 
meaning from code and data was genuinely refreshing and enjoyable. It felt like 
a far more honest way to learn than beginning from first principles.

My main contribution was extending the pipeline to pull FFIEC Call Report data 
for 2023–2025 and filling in the gaps across the codebase so that everything 
runs end-to-end. That second part was surprisingly enjoyable. It felt like 
solving a puzzle, where each missing piece I discovered helped me understand 
the system more clearly. Working out which variables belonged to which FFIEC 
schedules, why certain type conversions broke downstream, or how to parameterise 
a module so that it could serve both WRDS and FFIEC without breaking anything 
— each fix deepened my understanding of both the code and the underlying finance. 

%=============================================================================
\bibliographystyle{jpe}
\bibliography{bibliography}

\end{document}
